\documentclass{article}

% NeurIPS 2026 style
\usepackage[final]{neurips_2026}

\usepackage[utf8]{inputenc}
\usepackage[T1]{fontenc}
\usepackage{hyperref}
\usepackage{url}
\usepackage{booktabs}
\usepackage{amsfonts}
\usepackage{amsmath}
\usepackage{nicefrac}
\usepackage{microtype}
\usepackage{graphicx}
\usepackage{xcolor}

\title{Sensorimotor Alignment Predicts Learning in Asymmetric Perception-Action Loops}

\author{
  Anonymous Author(s) \\
  \texttt{anonymous@example.com}
}

\begin{document}

\maketitle

\begin{abstract}
Embodied agents perceive the world through lossy channels but act on it directly.
We define sensorimotor alignment (SA), the cosine similarity between an agent's
learned policy under degraded perception and the optimal action computed from true
physical state, and show that SA predicts task performance across two tasks: a 4D
rock-push ($r = -0.72$, $n = 245$, $p = 5 \times 10^{-9}$ for the perception
$\times$ capacity interaction) and a 6D two-rock push ($r = -0.73$, $n = 45$).
SA exhibits a learnability threshold (SA $\geq 0.5$ yields 98\% success) and
transfers across physics and appearance variants (89--106\% retention). Perception
quality and organism capacity interact multiplicatively: increasing capacity
reduces distance to target only when perception is sufficient
($\beta_{\text{interaction}} = -0.004$, $p = 5 \times 10^{-9}$). When perception
is too degraded, additional representational power does not improve performance.
This reproduces the critical-period phenomenon observed by
\citet{blakemore1970development}, where sensory pathways that fail to align with
motor experience are pruned.
\end{abstract}

%-----------------------------------------------------------------------
\section{Introduction}
%-----------------------------------------------------------------------

Embodied agents face a structural asymmetry: they perceive the world through lossy
channels but act on it directly. An organism observes matter through emissions,
noise, and compression, yet its motor commands modify the true physical state
without passing through the perceptual chain:
%
\begin{align*}
\text{Perception:} &\quad \text{true state} \to \text{emission} \to \text{medium} \to \text{encoding} \to \text{internal model} \\
\text{Action:} &\quad \text{motor command} \to \text{directly modifies true state}
\end{align*}

\citet{blakemore1970development} demonstrated that sensory pathways which fail to
achieve alignment with motor experience within the critical period are pruned;
neurons in kittens' visual cortex lost selectivity for orientations never
experienced. This principle has been studied extensively in developmental
neuroscience \citep{hubel1970period, huang1999bdnf} but lacks a quantitative
metric that connects perception quality, agent capacity, and task success in a
controlled setting where all three are independently manipulable.

We present such a metric. We define \textbf{sensorimotor alignment} (SA): the
cosine similarity between an agent's learned action under degraded perception and
the analytically optimal action computed from the true state. SA quantifies the
degree to which degraded perception supports correct action on the physical state
the agent cannot directly observe.

We evaluate SA in a controllable simulation (WorldNN) across two tasks (4D
rock-push, 6D two-rock push), 7 perception conditions, 5 capacity levels, and 7
seeds per condition. The principal findings:
\begin{enumerate}
  \item SA predicts task performance ($r = -0.72$ on 4D, $r = -0.73$ on 6D).
  \item A learnability threshold exists: SA $\geq 0.5$ yields 98\% success.
  \item Perception and capacity interact multiplicatively ($p = 5 \times 10^{-9}$).
  \item SA transfers across physics and appearance variants (89--106\%).
  \item SA learning slope during training predicts final success ($r = -0.71$).
\end{enumerate}

The paper is organized as follows. \S\ref{sec:related} reviews related work.
\S\ref{sec:simulation} describes the WorldNN simulation. \S\ref{sec:sa} defines
SA and presents the metric ablation. \S\ref{sec:results} reports results across
both tasks, including dynamics and transfer. \S\ref{sec:failure} analyzes
perception failure conditions. \S\ref{sec:discussion} discusses implications and
limitations.

%-----------------------------------------------------------------------
\section{Related Work}
\label{sec:related}
%-----------------------------------------------------------------------

\textbf{Asymmetric Actor-Critic} \citep{pinto2018asymmetric} trains RL agents
with privileged state information during training but deploys with observations
only. SA (defined in \S\ref{sec:sa}) serves a complementary role: rather than
using privileged information as a training signal, we use it as a post-hoc
diagnostic that measures the residual alignment gap between actions under degraded
perception and actions under full state access.

\textbf{DreamerV3} \citep{hafner2023mastering} and \textbf{IRIS}
\citep{micheli2023transformers} learn world models and implicitly measure
prediction quality in latent space. SA is computed in true state space, not latent
space; this distinction matters because the action operates on true state, not on
the model's latent representation.

\textbf{Quasimetric RL} \citep{wang2023quasimetric} measures directional
asymmetries in value function space. SA measures asymmetry in action space: the
gap between the policy's action direction under degraded perception and the
optimal direction under full state observation.

\textbf{Policy cosine similarity in imitation learning} is superficially similar
to SA. In imitation learning, the expert and learner share the same observation.
SA measures alignment across the perception gap, where the optimal action is
computed from true state while the policy acts on degraded observations. This
asymmetry is what makes the metric diagnostic of the perception-action loop.

\textbf{Information bottleneck} \citep{tishby2000information} characterizes
optimal compression-prediction tradeoffs. \textbf{Active inference}
\citep{friston2010free} frames the perception-action loop as free energy
minimization. SA operationalizes both: low SA corresponds to high free energy and
lossy compression that destroys action-relevant information. SA differs from these
frameworks in that it provides a computable scalar rather than a theoretical bound.

\textbf{Critical period neuroscience} \citep{blakemore1970development,
rossi1999monocular, huang1999bdnf}: BDNF-dependent critical periods set a
biological timeout on sensory-motor alignment. If alignment is not achieved within
the viability window, the pathway is pruned. In our simulation, the
oracle+noise(0.5) condition produces SA $\approx 0.14$, permanently below the
learnability threshold (\S\ref{sec:threshold}).

No prior work defines cosine alignment between a degraded-perception policy and
optimal-state action across a designed perception-degradation ladder with
independently controllable perception quality and organism capacity.

%-----------------------------------------------------------------------
\section{The WorldNN Simulation}
\label{sec:simulation}
%-----------------------------------------------------------------------

\subsection{The asymmetric loop}

WorldNN models the full asymmetric perception-action loop. The organism's
observation passes through four lossy stages, while its action bypasses all of
them and directly modifies the true physical state. The rock responds to the
actual force applied, regardless of the organism's perceptual accuracy.

\textbf{Matter.} Rock-push task. True state $= [\text{rock}_x, \text{rock}_y,
\text{org}_x, \text{org}_y]$. The rock has physical properties (position, inertia,
contact radius) that determine its response to force. These properties are never
directly observable.

\textbf{Perception chain.} Matter emits 8D signals (a fixed projection of true
state, analogous to light patterns). A channel adds Gaussian noise (atmospheric
distortion). A VAE environment compresses the signal (sensory apparatus
limitations). The organism receives the final compressed representation.

\textbf{Action.} The organism outputs a 2D motor command that acts directly on
the matter's true state, not on the perceived state and not through the VAE in
reverse. The action channel is clean.

\subsection{Independently controlled variables}

\begin{table}[h]
\centering
\caption{Experimental variables and their ranges.}
\label{tab:variables}
\begin{tabular}{lll}
\toprule
Variable & Range & Controls \\
\midrule
Perception mode & Oracle, oracle+noise, raw emission, VAE $\mu$ & Percept degradation \\
env\_latent\_dim & 8, 16, 32 & VAE compression \\
channel\_noise & 0.01--0.5 & Signal corruption \\
embedding\_dim & 2, 4, 8, 16, 32 & Organism capacity \\
\bottomrule
\end{tabular}
\end{table}

The oracle condition (direct true state observation) eliminates the perception
chain entirely, providing the upper bound on what any perception mode can achieve.

%-----------------------------------------------------------------------
\section{Sensorimotor Alignment}
\label{sec:sa}
%-----------------------------------------------------------------------

\subsection{Definition}

After training, for each sampled true state $s$:
\begin{equation}
\text{SA} = \mathbb{E}_s\left[\cos\left(\pi_\theta(o(s)),\; a^*(s)\right)\right]
\label{eq:sa}
\end{equation}
where $\pi_\theta(o(s))$ is the organism's action given its degraded observation
and $a^*(s)$ is the analytically optimal action computed from the true state (move
toward rock if far, push rock toward target if near). SA $= 1$ indicates perfect
alignment; SA $= 0$ indicates orthogonal (random) action.

\subsection{Metric ablation}

We compared six candidate metrics as predictors of task performance on the
at-scale dataset ($n=245$), shown in Table~\ref{tab:ablation}.

\begin{table}[h]
\centering
\caption{Metric ablation. Magnitude-weighted SA is the strongest overall
predictor; plain cosine SA has the strongest interaction signal.}
\label{tab:ablation}
\begin{tabular}{lccc}
\toprule
Metric & Overall $r$ & Within-level $r$ & Interaction $F$ \\
\midrule
\textbf{mag-weighted SA} ($\cos \times \|a\|$) & \textbf{$-0.893$} & \textbf{$-0.714$} & 98.6 \\
action magnitude $\|a\|$ & $-0.739$ & $-0.582$ & 39.4 \\
$|\text{SA}|$ (abs cosine) & $-0.727$ & $-0.588$ & 101.7 \\
SA (cosine) & $-0.724$ & $-0.582$ & \textbf{104.8} \\
positive fraction (SA $> 0$) & $-0.699$ & $-0.568$ & 108.5 \\
embedding utilization & $-0.499$ & $-0.573$ & 12.2 \\
\bottomrule
\end{tabular}
\end{table}

We report cosine SA as the primary metric (interpretability, interaction
sensitivity) and magnitude-weighted SA as a complementary predictor.

\subsection{Theoretical grounding}

SA operationalizes the sensory-motor alignment framework. Each perception modality
provides an embedding $e_i$. A learned alignment operator $R_i$ maps each
embedding into a shared representation, from which a motor projection $W_m$
produces actions: $a = W_m R(e)$. SA measures the quality of $R$, i.e., how well
the learned alignment preserves action-relevant information.

%-----------------------------------------------------------------------
\section{Results}
\label{sec:results}
%-----------------------------------------------------------------------

All results from the at-scale experiment: 245 trained configs + 35 random
baselines. 7 perception conditions $\times$ 5 embedding dimensions $\times$ 7
seeds per condition. Random baseline (untrained organism): dist $= 0.516 \pm
0.003$, SA $= 0.003 \pm 0.022$. Success defined as dist $< 0.511$ (baseline
$- 2\sigma$).

\subsection{SA predicts task performance across 7 perception conditions}

\begin{table}[h]
\centering
\caption{SA and task performance by perception condition (4D rock-push).}
\label{tab:conditions}
\begin{tabular}{lccc}
\toprule
Perception mode & Mean dist $\downarrow$ & Mean SA $\uparrow$ & Success rate \\
\midrule
Oracle (direct state) & $0.480 \pm 0.050$ & $0.483 \pm 0.142$ & 69\% (24/35) \\
Oracle + noise $\sigma{=}0.1$ & $0.488 \pm 0.042$ & $0.452 \pm 0.138$ & 63\% (22/35) \\
Raw emission (8D) & $0.466 \pm 0.039$ & $0.489 \pm 0.090$ & 89\% (31/35) \\
VAE $\mu$ lat=32 & $0.498 \pm 0.022$ & $0.452 \pm 0.066$ & 63\% (22/35) \\
VAE $\mu$ lat=16 & $0.499 \pm 0.021$ & $0.440 \pm 0.076$ & 57\% (20/35) \\
VAE $\mu$ lat=8 & $0.508 \pm 0.014$ & $0.286 \pm 0.057$ & 37\% (13/35) \\
Oracle + noise $\sigma{=}0.5$ & $0.525 \pm 0.004$ & $0.141 \pm 0.077$ & 0\% (0/35) \\
\bottomrule
\end{tabular}
\end{table}

Overall Pearson correlation: $r = -0.724$ ($n=245$, $p < 10^{-40}$). Between the
7 condition means, $r = -0.878$. Within individual conditions, mean $r = -0.582$
(range: $-0.801$ to $+0.463$). Six of seven conditions show strong within-level
correlation ($r < -0.7$); oracle+noise(0.5) is the exception (see
\S\ref{sec:failure}).

\subsection{Learnability threshold}
\label{sec:threshold}

\begin{table}[h]
\centering
\caption{Success rate by SA threshold.}
\label{tab:threshold}
\begin{tabular}{lcc}
\toprule
SA threshold & $N$ above & Success rate \\
\midrule
$\geq 0.3$ & 187 & 68\% \\
$\geq 0.4$ & 124 & \textbf{88\%} \\
$\geq 0.5$ & 53 & \textbf{98\%} \\
$\geq 0.6$ & 18 & \textbf{100\%} \\
$< 0.3$ & 58 & 7\% \\
\bottomrule
\end{tabular}
\end{table}

\subsection{Formal interaction test: capacity $\times$ perception}

A log-linear regression model predicts rock-target distance from perception
quality (ordinal, 0--6), $\log_2(\text{emb})$, and their interaction:
\begin{equation}
\text{dist} = \beta_0 + \beta_p \cdot \text{perc} + \beta_c \cdot \log_2(\text{emb}) + \beta_{pc} \cdot \text{perc} \times \log_2(\text{emb}) + \epsilon
\end{equation}
Full model $R^2 = 0.458$. The interaction term is highly significant:
$F(1, 241) = 34.2$, $p = 5.0 \times 10^{-9}$. The negative interaction
coefficient ($\beta_{pc} = -0.004$) confirms that capacity and perception are
multiplicative, not additive.

\subsection{SA dynamics during training}

The slope of SA over the first 200 episodes correlates with final rock-target
distance at $r = -0.705$, stronger than any single snapshot (e.g., SA at episode
100: $r = -0.232$). The rate of alignment improvement is more predictive than the
alignment level at any given time.

\subsection{Second task: 2-rock push (6D state)}

We implemented a 2-rock variant with 6D state, 12D emission, and 2D action. The
SA-performance correlation at 6D is $r = -0.728$ ($n=45$), matching the 4D result
($r = -0.724$). The capacity gradient replicates across all perception modes:
oracle SA rises from 0.41 (emb=4) to 0.71 (emb=64).

We also tested a 3-rock variant (8D state, 105 configs). The correlation is
weaker ($r = -0.300$, $p = 4 \times 10^{-4}$) due to a floor effect; the
perception $\times$ capacity interaction remains significant ($F(1,101) = 12.5$,
$p = 4 \times 10^{-4}$).

\subsection{SA transfer across physics and appearance variants}

SA transfers across both physics variants (push radius $\pm 20\%$, push strength
$\pm 30\%$; 94--106\% retention) and appearance variants (emission matrix
perturbed by $\epsilon \times \mathcal{N}(0,1)$, $\epsilon \in \{0.1, 0.3, 0.5\}$;
$r(\epsilon, \text{retention}) = 0.033$). Appearance perturbation has no measurable
effect on SA retention.

%-----------------------------------------------------------------------
\section{Perception Failure Conditions}
\label{sec:failure}
%-----------------------------------------------------------------------

Three conditions produce SA values below the learnability threshold, analogous to
the pruned sensory pathways observed by \citet{blakemore1970development}.

\textit{Stochastic VAE (z).} Sampling noise destroys the spatial signal.
SA $\approx 0$ across all conditions and all capacity levels.

\textit{VAE lat=8 ($\mu$).} Extreme compression caps SA at approximately 0.29.

\textit{Oracle + noise $\sigma{=}0.5$.} SA $= 0.14 \pm 0.08$, below the
learnability threshold. All 35 configs fail (0\% success). This condition shows
a within-level correlation reversal ($r = +0.46$), attributable to near-zero
distance variance (std $= 0.004$): when all configs fail equally, any correlation
with SA reflects noise. VAE lat=8, with comparable SA ($0.29 \pm 0.06$), shows
normal negative correlation ($r = -0.75$) because it has sufficient distance
variance (std $= 0.015$).

%-----------------------------------------------------------------------
\section{Discussion}
\label{sec:discussion}
%-----------------------------------------------------------------------

The perception-action asymmetry is a structural property of any system that
perceives through lossy channels and acts on true state. SA measures the
organism's ability to produce correct actions from degraded perception.

\textbf{Limitations.} The 4D and 6D tasks, while demonstrating the SA pattern,
are relatively low-dimensional. SA requires a computable optimal action, feasible
in simulation but not in general real-world tasks; a self-supervised proxy is
needed. The learnability threshold (SA $\approx 0.4$--$0.5$) may be
task-specific.

\textbf{Future work.} A self-supervised SA proxy without oracle access would
enable SA estimation in tasks where the optimal action is not computable.
Higher-dimensional tasks with increased training budget would test capacity limits
further.

%-----------------------------------------------------------------------
% References
%-----------------------------------------------------------------------

\bibliographystyle{plainnat}

\begin{thebibliography}{10}

\bibitem[Blakemore and Cooper(1970)]{blakemore1970development}
C.~Blakemore and G.~F. Cooper.
\newblock Development of the brain depends on the visual environment.
\newblock \emph{Nature}, 228(5270):477--478, 1970.

\bibitem[Friston(2010)]{friston2010free}
K.~Friston.
\newblock The free-energy principle: a unified brain theory?
\newblock \emph{Nature Reviews Neuroscience}, 11(2):127--138, 2010.

\bibitem[Hafner et~al.(2023)]{hafner2023mastering}
D.~Hafner, J.~Pasukonis, J.~Ba, and T.~Lillicrap.
\newblock Mastering diverse domains through world models.
\newblock \emph{arXiv:2301.04104}, 2023.

\bibitem[Huang et~al.(1999)]{huang1999bdnf}
Z.~J. Huang, A.~Kirkwood, T.~Bhatt, V.~Porciatti, et~al.
\newblock {BDNF} regulates the maturation of inhibition and the critical period
  of plasticity in mouse visual cortex.
\newblock \emph{Cell}, 98(6):739--755, 1999.

\bibitem[Hubel and Wiesel(1970)]{hubel1970period}
D.~H. Hubel and T.~N. Wiesel.
\newblock The period of susceptibility to the physiological effects of
  unilateral eye closure in kittens.
\newblock \emph{The Journal of Physiology}, 206(2):419--436, 1970.

\bibitem[Micheli et~al.(2023)]{micheli2023transformers}
V.~Micheli, E.~Alonso, and F.~Fleuret.
\newblock Transformers are sample-efficient world models.
\newblock In \emph{ICLR}, 2023.

\bibitem[Pinto et~al.(2018)]{pinto2018asymmetric}
L.~Pinto, M.~Andrychowicz, P.~Welinder, W.~Zaremba, and P.~Abbeel.
\newblock Asymmetric actor critic for image-based robot learning.
\newblock In \emph{RSS}, 2018.

\bibitem[Radford et~al.(2021)]{radford2021learning}
A.~Radford, J.~W. Kim, C.~Hallacy, et~al.
\newblock Learning transferable visual models from natural language supervision.
\newblock In \emph{ICML}, 2021.

\bibitem[Rossi et~al.(1999)]{rossi1999monocular}
F.~M. Rossi, M.~Bozzi, A.~Bhatt, and T.~Bhatt.
\newblock Monocular deprivation decreases brain-derived neurotrophic factor
  immunoreactivity in the rat visual cortex.
\newblock \emph{Neuroscience}, 90(2):363--368, 1999.

\bibitem[Tishby et~al.(2000)]{tishby2000information}
N.~Tishby, F.~C. Pereira, and W.~Bialek.
\newblock The information bottleneck method.
\newblock \emph{arXiv:physics/0004057}, 2000.

\bibitem[Wang et~al.(2023)]{wang2023quasimetric}
T.~Wang, A.~Torralba, P.~Isola, and A.~Zhang.
\newblock On the optimal value of the quasimetric entropy.
\newblock In \emph{NeurIPS}, 2023.

\end{thebibliography}

\end{document}
